\hypertarget{appendix-bc44f7b42f42}{}
\section{Unit roundoff}\label{appendix-bc44f7b42f42-unit-roundoff}

\textbf{Curriculum depth:} Developed · \textbf{First seen:} C004 · \textbf{Latest development:} C005

\subsection{Definition and mechanism}\label{appendix-bc44f7b42f42-definition-and-mechanism}

Under nearest binary rounding the conventional relative bound for normal-range results is u=2\^{}-p, half upward epsilon. For binary32 this is 2\^{}-24. The bound has range/rounding assumptions.

Repeated-rounding bounds use \(\gamma_n=nu/(1-nu)\) when \(nu<1\), under the stated rounding and range assumptions. Summation error also depends on operand magnitudes; this is not a universal relative-error guarantee.

\subsection{Prerequisites}\label{appendix-bc44f7b42f42-prerequisites}

\begin{itemize}
\tightlist
\item
  \hyperref[appendix-14bab4cab9d4]{ULP and representable spacing} --- Adjacent representable points define rounding candidates, errors and step distances.
\item
  \hyperref[appendix-614afe1de592]{Rounding to nearest ties to even} --- Midpoint parity determines discarded-tail decisions and the normal relative-error bound.
\end{itemize}

\subsection{Used by}\label{appendix-bc44f7b42f42-used-by}

\begin{itemize}
\tightlist
\item
  \hyperref[appendix-3ae74e89b4e6]{Absolute and relative error} --- The local normal-rounding bound supplies the error-analysis scale.
\item
  \hyperref[appendix-f82e270791c8]{Accumulated rounding error} --- The local normal-rounding bound supplies the error-analysis scale.
\end{itemize}

Connect later dependent concepts here as their lessons arrive.

\begin{itemize}
\tightlist
\item
  \href{../02-Foundations-II/main.pdf\#nameddest=appendix-d580b422240f}{Stable probability normalization and log-sum-exp (02-Foundations-II, appendix-d580b422240f)}
\item
  \href{../02-Foundations-II/main.pdf\#nameddest=appendix-b9afba7f63e5}{Stable quadratic root evaluation (02-Foundations-II, appendix-b9afba7f63e5)} \#\# Recall and next depth
\end{itemize}

Reconstruct the concrete example in the definition and identify its validity assumptions.

\textbf{Pending:} Subnormal absolute-error models and target-specific arithmetic require separate treatment; do not extend the normal relative bound to arbitrarily tiny results.

\subsection{Source and navigation}\label{appendix-bc44f7b42f42-source-and-navigation}

\hyperref[01-foundations-i-ch001-when-arithmetic-meets-uncertainty]{M005 development}

\hyperref[01-foundations-i-ch001-a-stretching-ruler-for-real-valued-quantities]{M004 teaching chapter} · \hyperref[appendix-fd090d982a15-m004---parts-61-76]{Per-part definitions} · \href{../02-Foundations-II/main.pdf\#nameddest=appendix-4ad47bdd269a}{Dependency Map (02-Foundations-II, appendix-4ad47bdd269a)}
